\section{论文按时完成的可能性}

从当前进度来看，本课题已经具备按时完成论文的现实基础。其主要原因在于：研究问题已经明确为面向高代价蛋白质设计工作流的自适应工具链规划算法；核心算法主链已经形成较完整的设计、实现和验证基础；系统运行骨架、恢复链、等待态机制、日志快照和实验入口也已较为成熟。换言之，当前剩余工作更多是围绕实验收口、量化统计和论文表达展开，而不是重新发起大规模系统重构。

与此同时，后续工作虽然仍有压力，但边界已经比较清楚。横向对比实验、工作量统计补充、统一图表整理和文字最终收束，均建立在现有系统与现有算法主链之上推进，属于“补齐和整合”性质，而不属于“从零开始搭建”的工作。这意味着只要后续继续保持围绕核心算法的统一叙事，并按既定顺序完成实验和统计补充，论文整体收口仍处于可控范围之内。

因此，可以认为本课题已经具备按时完成论文工作的较高可能性。后续的关键不在于继续扩展系统边界，而在于把已经形成的核心方法、系统支撑和实验入口整理为更完整、更可信、更具量化支撑的论文表达。
